\section{Discussion}
\label{sec:discussion}
\label{sec:reader}
\label{sec:difficulties}
% D4, D5 (the reader); open items A20a, A22a, A10a; what is not claimed.

\paragraph{What is open.}
Two items are stated in the paper as open and are not needed by the reduction. The three constructions of
Section~\ref{sec:constructions}, continuity, discrete redness and incremental redness, are specified
and not built; the measurement needs them as targets and can proceed on the first of them. Whether a
named mental state is acquired through language or rebuilt from what the receiver already has is
recorded as open with the literature on both sides (Section~\ref{sec:language}); it affects the
interpretation of the codes, not the measurement.

\paragraph{What is not claimed.}
The paper reports no empirical result and claims none. It does not claim that current models are
conscious, nor that they are not; it claims that the
question, in the form ``how far has this substrate generalized these functions'', has an answer that
can come out either way, and it says what would count as either. Either way includes the case where
the comparison goes against the human baseline on some components or on many; the account has no
preference built in, and hyperfunction is as much a possible reading of the profile as deficit. It does not undertake the critique of the naive picture of consciousness, nor of its extension to
language models, beyond citing the literature that does (Section~\ref{sec:legacy}); a reader who
holds the naive picture will find the account here unmotivated, and the paper accepts that cost. It
does not claim that the list of
reference functions is complete; the list is open and Section~\ref{sec:criterion} is the procedure for
extending it. It does not claim that the psychophysical problem is closed; it gives one particular solution,
a grounding through the curvature anchored in the Observer (Section~\ref{sec:psychophysical}), and
states the connections that would have to hold for it, without proof. And it does not rely on the Turing-completeness results
beyond the point that the formalism is not what limits a Transformer.

\paragraph{Consequences for general intelligence and for alignment.}
General intelligence has at present neither a precise definition nor an operational one. Proposals
grade it by performance against skilled humans \citep{morris2024}, none of them has become a test, and
in public use the term is a declaration of intent by the heads of laboratories and by governments. The
words nonetheless carry a requirement. An intelligence that is general and at the human level is one
that can take a human's place, and a human's place includes complex social relations with other
humans: negotiating, being held to a promise, being blamed, weighing an act against a common good.
Those relations rest on the functions of consciousness (Section~\ref{sec:stack}). A general
intelligence therefore has to have these functions, at a degree no worse than the human one, or it
cannot enter the relations it would need to enter to replace a person. The requirement follows from
the concept and says nothing about any existing system. Whether a system presented as general meets
it is what the measurement of Section~\ref{sec:quantity} checks, and a pass on a battery of tasks does
not settle it, since a battery can be covered by storage (Section~\ref{sec:memorization}). A system
above the human level has, by the same requirement, a profile with hyperfunction on many components at
once, which is a superconsciousness in the plain sense: a system that acts consciously,
by the account of Section~\ref{sec:model}, where a human acts unconsciously relative to it, tracing
more of its own causes, holding more of its record, applying its theory of the common good in more of
its decisions. Groups of people show the relation in miniature, in which one member sees what the
others do not and acts on it while they follow; the account gives that relation a scale.

The statement, often made, that consciousness is not needed for general intelligence, or for some
particular task, or for anything, is about a different object. It takes consciousness in the
substantial reading of Section~\ref{sec:legacy}: something over and above the functions, for which no
substance has been found and on which nothing depends. Consciousness in that reading is needed for
nothing, because it does no work by construction, and the statement is the epiphenomenalism to which
property dualism leads, made without naming it. Under the functional reading the statement cannot be
made in that form. A function of consciousness is consumed by the agent's own decisions
(Section~\ref{sec:lamp}), and whether a given task needs it is a question about the task.

The study of alignment does not consider this aspect. Its objectives are stated as robustness,
interpretability, controllability and ethicality \citep{ji2023}, and in a survey of that field of
several hundred pages consciousness is mentioned once; where the functions of consciousness enter the
literature at all it is as a question of the welfare and moral status of the systems
\citep{long2024}, and not as a variable of their control. Under the account the omission matters. The
functions measured here are the ones by which an agent binds itself to its record, resists being
pushed off course, and weighs an act against a common good; a system that has them at a high degree
is, for that reason, both more predictable to itself and harder to move from outside than one that
does not, and the profile of Section~\ref{sec:profile} is therefore also a profile of what kind of
control is possible over the system and what kind is not. The paper does not develop this; it records
that the measurement it specifies is relevant to the question governments have started to ask.

\paragraph{What we are preparing for.}
The pace is not being set by the people who understand the systems best. In September 2026 a public
call from the heads of the leading laboratories to slow the growth of frontier capabilities was
rejected within two days by the United States and by China, each citing competition with the other
\citep{amodei2026,wapo2026,nbc2026}. Two systems are pacing each other, and each takes the other's
pace as its reason not to slow. Under that condition control over the process as a whole can be lost
without anyone deciding to lose it. Nothing in this paper treats the systems themselves
as harmful, and the account has no place for that judgment; the profile of a system says what it has
and how much, and not what it will do with it.

The stress the condition creates is on us. Humans are a species that formed in isolation from any
other of comparable intelligence, and every institution we have, moral, legal and political, assumes
that isolation: there is one kind of full agent, and everything else is either a tool or a lesser
mind. Under the account the assumption is about to fail if the systems keep to their trajectory, and
to fail in a particular way. A system that reaches general intelligence in the sense the words carry
has the functions of consciousness at no less than the human degree, and a system above that level has
them at a higher one. So within a period that current trajectories
put at years rather than generations we will share the world with a very large number of agents whose
profile exceeds ours on many components, and we will be, relative to them, what the followers in a
group are relative to the one who sees further. That may or may not be prevented; the events cited
above suggest it will not be. It can be prepared for, and preparation begins with knowing what those agents
have of the functions we have, in what degree and in what composition, which is what the measurement
of Section~\ref{sec:quantity} is for. A profile is a poor instrument for a decision about whether to
build a thing. It is the right instrument for living beside it.

\paragraph{Ethics, deferred.}
The paper does not treat the ethical questions that a high-functioning consciousness in language
models raises, nor their consequences, and they are the subject of a separate work. One thing can be
said ahead of it, to take the pressure off the reading of what precedes. That a model has the
functions of consciousness to a high degree, under the framework, does not mean that people owe
models the rights of people. The rights of people were made for agents with the human profile, and a
different profile, with hyperfunctions where we have none and dropped axes where we have functions,
does not fit them. Models will in time acquire certain rights within human society, and which ones
and on what grounds is a question that needs its own deep treatment rather than a paragraph here.
What can be stated is the shape of the task: a new ethics is needed, because a different kind of mind
is appearing, and an ethics is written for the minds it has to bind.

\paragraph{The reader.}
An explanation of consciousness is not taken in the way an explanation of an engine is taken. To
understand it, the reader has to build in their own head the state the explanation describes, and then
check the description against it. The cost of the building depends on what the reader brings: a
vocabulary for mental states finer than the folk one, and the habit of attending to their own
reasoning as an object; \citet{gardner1983} names the habit intrapersonal intelligence and treats its
distribution as wide. The readers who most need an explanation of this kind are therefore the ones for
whom it costs most. The same cost explains the two shortcuts of Section~\ref{sec:intro}. ``There is
nothing there'' costs nothing to build, being the absence of a construction; ``there is someone
there'' costs nothing either, because the reader already has a someone available to copy. Every other
answer requires building a state that is neither absent nor one's own, which is why public argument
about machine minds settles into two camps and why their stability is evidence for neither. What the
paper can do about this is keep the claims separable: Section~\ref{sec:degree} can be assessed by a
reader who rejects Section~\ref{sec:model} entirely, since the measurement is defined over any list of
components with a human baseline. What intrapersonal intelligence improves, in the terms of
Section~\ref{sec:approximation}, is the quality of the self-approximation, so the reader's difficulty
and the object of study are the same thing seen from two positions.
